%% =====================================================================
%%  Unit 20 — Happy birthday, Léa!  (Faire le point)
%%  Pistes : 20.01–20.04 dialogues · 20.05–20.06 vocabulaire · 20.07–20.08 grammaire
%%           20.09 prononciation · 20.10–20.13 exercices (textes en annexe)
%% =====================================================================

\newcommand{\otIllusUnXX}{\otIllusDuo{cake/104/86/15/0.10,party/190/92/20/0.09,wish/122/14/13/0.08,
  candles/150/8/11/0.08,cheers/196/22/12/0.07,friends/100/50/12/0.07}{emma}{Surprise!}{lea}{Oh, thank you!}}

% Carte d'anniversaire (exercice 4)
\newcommand{\otCarteAnniv}[1]{%
  \begin{tcolorbox}[enhanced,width=124mm,center,colback=white,colframe=otViolet,boxrule=0.8pt,arc=2.5mm,
    left=8mm,right=8mm,top=12mm,bottom=6mm,fontupper=\fontsize{10.6}{14.4}\selectfont\itshape,
    overlay={%
      \foreach \x/\y/\c/\r in {6/-6/otOrange/1.6,14/-9/otTeal/1.1,22/-5/otRed/1.3,104/-6/otGreen/1.5,
        112/-10/otViolet/1.1,118/-5/otOrange/1.2,28/-10/otViolet/0.9,96/-11/otTeal/0.9}{%
        \fill[\c] ([xshift=\x mm,yshift=\y mm]frame.north west) circle (\r mm);}
      \node[font=\otBold\fontsize{17}{20}\selectfont,text=otViolet] at ([yshift=-7mm]frame.north)
        {\enb{Happy Birthday!}};}]
  #1
  \end{tcolorbox}}

\unite[Happy birthday,\\Léa!]{20}{Happy birthday, Léa!}{Faire le point}{otIllusUnXX}

\situation{Épisode 20}{Le 15 mars, l'anniversaire de Léa. Emma lui a donné rendez-vous au café de Sam
« pour un simple verre »… mais toute la bande l'attend dans le noir. Six mois après son arrivée à New
York, c'est l'heure de faire le point !}

\begin{objectifs}
\item parler de vous au passé, au présent et au futur ;
\item choisir le bon temps grâce aux mots repères (\en{yesterday}, \en{now}, \en{next year}…) ;
\item poser toutes sortes de questions en trouvant le bon auxiliaire ;
\item souhaiter un anniversaire, féliciter, remercier ;
\item respecter le rythme de la phrase anglaise.
\end{objectifs}

\auprogramme
  {Parler de soi au passé, au présent et au futur}
  {Révision : les temps et les questions}
  {Fêtes, cadeaux, souhaits : 21 expressions}
  {Le rythme de la phrase anglaise}
  {Anniversaires et fêtes}

\semaine{casque/otNavy/Dialogues/Écouter et lire les deux scènes,
  bulle/otGreenDark/Vocabulaire/Apprendre les 21 expressions,
  blocs/otViolet/Grammaire/Les temps et les questions,
  sons/otOrangeDark/Sons et culture/Le rythme et les fêtes,
  crayon/otNavy/Exercices/12 exercices corrigés,
  coche/otGreenDark/Bilan/L'essentiel et le test express}

\clearpage
%% ---------------------------------------------------------------------
\rubrique{ecouter}{Surprise !}

\begin{dialogue}{1}{Surprise!}{Vendredi 15 mars, 19 heures. Léa pousse la porte du café de Sam. La
  lumière s'allume d'un coup.}{20.01}{20.02}
\dit{emma}{Surprise! \hl{Happy birthday}, Léa!\note{1}}{Surprise ! Joyeux anniversaire, Léa !}
\dit{lea}{Oh my goodness! \hl{Did you plan} all this?\note{2}}{Oh là là ! C'est vous qui avez organisé tout ça ?}
\dit{jack}{Yes, we did! Emma \hl{had} the idea, and Sam \hl{made} the cake.}{Oui ! Emma a eu l'idée, et Sam a fait le gâteau.}
\dit{sam}{It's a chocolate cake. I \hl{made} it this morning!}{C'est un gâteau au chocolat. Je l'ai fait ce matin !}
\dit{lopez}{\hl{How old are you} today, Léa?\note{3}}{Quel âge avez-vous aujourd'hui, Léa ?}
\dit{lea}{I'm twenty-eight! Thank you so much, everybody!}{J'ai vingt-huit ans ! Merci beaucoup, tout le monde !}
\dit{maya}{Come on, \hl{make a wish} and blow out the candles!\note{4}}{Allez, fais un vœu et souffle les bougies !}
\dit{jack}{So, \hl{what did you wish for}?\note{5}}{Alors, qu'est-ce que tu as souhaité ?}
\dit{lea}{I can't tell you. It's a secret!}{Je ne peux pas te le dire. C'est un secret !}
\end{dialogue}

\begin{comprendre}
\pointnote{1} \en{\textbf{Happy birthday!}} « Joyeux anniversaire ! ». Pour Noël, on dit \en{Merry
  Christmas!} ; pour la nouvelle année, \en{Happy New Year!}
\pointnote{2} \en{\textbf{Did you plan…?}} Une question au prétérit (unité 18) : \en{did} + base
  verbale. \en{plan} = organiser, prévoir.
\pointnote{3} \en{\textbf{How old are you?}} L'âge se dit avec \en{be} (unité 3) : \en{I'm
  twenty-eight}, jamais \en{I have twenty-eight}.
\pointnote{4} \en{\textbf{make a wish}} « faire un vœu », avant de souffler les bougies (\en{blow out
  the candles}).
\pointnote{5} \en{\textbf{What did you wish for?}} En anglais, la préposition se place souvent à la fin
  de la question : \en{What did you wish for?}, \en{Who are you talking to?}
\end{comprendre}

\begin{dialogue}{2}{Six months in New York}{Un peu plus tard, sur la petite terrasse du café. Emma et
  Léa prennent l'air.}{20.03}{20.04}
\dit{emma}{So, how do you feel? Six months in New York!}{Alors, comment tu te sens ? Six mois à New York !}
\dit{lea}{\hl{Six months ago}, I \hl{didn't know} anybody here. \hl{Now} I \hl{have} wonderful friends.}{Il y a six mois, je ne connaissais personne ici. Maintenant, j'ai des amis formidables.}
\dit{emma}{And your English is great now!}{Et ton anglais est super, maintenant !}
\dit{lea}{Thanks! In September, I \hl{was} so nervous. Now I \hl{speak} English every day.\note{6}}{Merci ! En septembre, j'étais si stressée. Maintenant, je parle anglais tous les jours.}
\dit{emma}{I'm proud of you! So, \hl{what are you going to do} next?}{Je suis fière de toi ! Alors, qu'est-ce que tu vas faire maintenant ?}
\dit{lea}{\hl{I'm going to take} an English class. And \hl{next summer}, I'm going to drive on Route 66 with Jack and Sam!}{Je vais suivre un cours d'anglais. Et l'été prochain, je vais faire la Route 66 en voiture avec Jack et Sam !}
\dit{emma}{Really? Can I come too?}{C'est vrai ? Je peux venir aussi ?}
\dit{lea}{Of course! \hl{Let's go} together!}{Bien sûr ! Partons tous ensemble !}
\end{dialogue}

\begin{comprendre}
\pointnote{6} \en{\textbf{I was… Now I speak…}} Léa résume son parcours avec tous les temps de ce
  livre : le passé (\en{I didn't know}, \en{I was}), le présent (\en{I have}, \en{I speak}) et le futur
  (\en{I'm going to take}). C'est l'objectif de cette dernière unité !
\end{comprendre}

%% ---------------------------------------------------------------------
\rubrique{mots}{Les mots de l'unité}

Voici les 21 expressions de l'unité : de quoi fêter un anniversaire… et faire le bilan de votre
propre parcours.

\begin{vocabulaire}{Fêtes et cadeaux}{20.05}
\categorie{La fête}
\mot{a birthday party}{ə ˈbɝːθdeɪ ˌpɑːrti}{une fête d'anniversaire}
\mot{a surprise party}{ə sɚˈpraɪz ˌpɑːrti}{une fête surprise}
\mot{a cake}{ə ˈkeɪk}{un gâteau}
\mot{a candle}{ə ˈkændəl}{une bougie}
\mot{make a wish}{ˌmeɪk ə ˈwɪʃ}{faire un vœu}
\mot{blow out the candles}{ˌbloʊ ˌaʊt ðə ˈkændəlz}{souffler les bougies}
\mot{celebrate}{ˈseləbreɪt}{fêter, célébrer}
\categorie{Cadeaux et souhaits}
\mot{a present}{ə ˈprezənt}{un cadeau}
\mot{a card}{ə ˈkɑːrd}{une carte}
\mot{Happy birthday!}{ˌhæpi ˈbɝːθdeɪ}{Joyeux anniversaire !}
\mot{Congratulations!}{kənˌɡrætʃəˈleɪʃənz}{Félicitations !}
\mot{Cheers!}{tʃɪrz}{Santé ! (pour trinquer)}
\end{vocabulaire}

\begin{vocabulaire}{Faire le bilan}{20.06}
\categorie{Parler de soi}
\mot{six months ago}{ˌsɪks ˈmʌnθs əˈɡoʊ}{il y a six mois}
\mot{nervous}{ˈnɝːvəs}{stressé, nerveux}
\mot{proud}{praʊd}{fier}
\mot[to improve]{improve}{ɪmˈpruːv}{progresser, s'améliorer}
\categorie{Phrases utiles}
\mot{I'm proud of you!}{aɪm ˈpraʊd əv juː}{Je suis fier de toi !}
\mot{Time flies!}{ˌtaɪm ˈflaɪz}{Le temps passe vite !}
\mot{Thank you for everything!}{ˈθæŋk juː fɚ ˈevriθɪŋ}{Merci pour tout !}
\mot{Let's celebrate!}{ˌlets ˈseləbreɪt}{Fêtons ça !}
\mot{Best wishes!}{ˌbest ˈwɪʃɪz}{Tous mes vœux !}
\end{vocabulaire}

%% ---------------------------------------------------------------------
\rubrique{grammaire}{Révision : les temps et les questions}

\otIntertitre[otViolet]{1 · Les cinq temps de ce livre}
\begin{center}\otTableau
\begin{tabular}{@{}l@{\hspace{6mm}}>{\enOn}l<{\enOff}@{\hspace{6mm}}>{\enOn}l<{\enOff}@{\hspace{6mm}}l@{}}
\otTH{Temps} & \otTH{Exemple} & \otTH{Mots repères} & \otTH{Unités}\\\midrule
présent simple (habitudes)     & I \textbf{work} every day.                 & every day, usually, on Mondays & 6, 7\\
présent en \en{-ing} (maintenant) & I'\textbf{m working} now.                 & now, right now, Look!          & 13\\
futur proche (projets)         & I'\textbf{m going to work} tomorrow.        & tomorrow, next week, soon      & 16\\
passé de \en{be}               & I \textbf{was} at work.                     & yesterday, last week, ago      & 17\\
prétérit                       & I \textbf{worked} / I \textbf{went} yesterday. & yesterday, last year, ago   & 18, 19\\\bottomrule
\end{tabular}\end{center}

\blocaudio{20.07}{Écoutez et répétez}{Le même verbe à tous les temps.}

\otIntertitre[otViolet]{2 · Les questions : trouvez l'auxiliaire}
\begin{center}\otTableau
\begin{tabular}{@{}>{\enOn\otSemi\color{otViolet}}l<{\enOff}@{\hspace{8mm}}>{\enOn}l<{\enOff}@{\hspace{8mm}}>{\enOn}l<{\enOff}@{}}
\otTH{Auxiliaire} & \otTH{Question} & \otTH{Réponse courte}\\\midrule
do / does   & \textbf{Do} you speak English? · \textbf{Does} she work here?   & Yes, I do. · No, she doesn't.\\
am / is / are & \textbf{Are} you working? · \textbf{Is} he going to come?    & Yes, I am. · No, he isn't.\\
can         & \textbf{Can} you swim?                                         & Yes, I can. · No, I can't.\\
was / were  & \textbf{Was} it good? · \textbf{Were} they happy?              & Yes, it was. · No, they weren't.\\
did         & \textbf{Did} you like it? · Where \textbf{did} you go?         & Yes, I did. · No, I didn't.\\\bottomrule
\end{tabular}\end{center}

\begin{regle}[La méthode en trois étapes]
\begin{etapes}
\item Trouvez le temps grâce aux mots repères (\en{yesterday} → prétérit ; \en{now} → présent en
  \en{-ing}…).
\item Placez l'auxiliaire devant le sujet ; ajoutez le mot interrogatif en tête si besoin
  (\en{Where}, \en{What}, \en{When}, \en{How}…).
\item Dans la réponse courte, reprenez le même auxiliaire : \en{Did you…? — Yes, I did.}
\end{etapes}
\end{regle}

\blocaudio{20.08}{Écoutez et répétez}{Cinq questions avec un auxiliaire différent, et leur réponse.}

\begin{attention}[Les pièges qui restent]
\fauxjuste{I have twenty-eight years.}{I'm twenty-eight.}
\fauxjuste{She don't like cake.}{She doesn't like cake.}
\fauxjuste{Did you went to the party?}{Did you go to the party?}
\fauxjuste{I'm going travel next year.}{I'm going to travel next year.}
\fauxjuste{Yesterday I am at home.}{Yesterday I was at home.}
\end{attention}

%% ---------------------------------------------------------------------
\rubrique{sons}{Le rythme de la phrase anglaise}

Le français donne à peu près la même durée à chaque syllabe. L'anglais, lui, avance au rythme des
\textbf{syllabes accentuées} : ce sont elles qui battent la mesure, et les petits mots entre elles
(\en{and}, \en{a}, \en{the}, \en{to}, \en{of}…) se compriment pour tenir dans le temps. Tapez dans vos
mains sur chaque mot en gras : les trois premières lignes durent presque le même temps !

\begin{phrases}{20.09}{Écoutez et répétez en tapant dans vos mains sur chaque syllabe en gras.}
\phrase[One, two, three, four.]{\textbf{One}, \textbf{two}, \textbf{three}, \textbf{four}.}{4 temps}
\phrase[One and two and three and four.]{\textbf{One} and \textbf{two} and \textbf{three} and \textbf{four}.}{4 temps : même durée !}
\phrase[One and a two and a three and a four.]{\textbf{One} and a \textbf{two} and a \textbf{three} and a \textbf{four}.}{4 temps : même durée !}
\phrase[Happy birthday, Léa!]{\textbf{Hap}py \textbf{birth}day, \textbf{Lé}a!}{3 temps}
\phrase[I'm going to see my friends.]{I'm \textbf{go}ing to \textbf{see} my \textbf{friends}.}{3 temps}
\phrase[I went to the beach with my friends.]{I \textbf{went} to the \textbf{beach} with my \textbf{friends}.}{3 temps}
\end{phrases}

%% ---------------------------------------------------------------------
\rubrique{culture}{Anniversaires et fêtes}

\begin{culture}[Un anniversaire à l'américaine]
On chante \en{Happy Birthday to You}, l'une des chansons les plus connues au monde. Le héros du jour
fait un vœu en silence (\en{make a wish}) et souffle ses bougies d'un coup : le vœu doit rester secret
pour se réaliser ! Certains âges comptent double : \en{Sweet Sixteen} (16 ans, souvent l'âge du permis
de conduire) et 21 ans, l'âge légal pour boire de l'alcool aux États-Unis.
\end{culture}

\begin{paire}
\begin{culture}[Offrir un cadeau]
Aux États-Unis, on ouvre son cadeau tout de suite, devant celui qui l'offre, et l'on s'extasie :
\en{Oh, I love it! You shouldn't have!} La carte (\en{a card}), avec un petit mot écrit à la main,
compte souvent autant que le cadeau.
\end{culture}
\begin{culture}[Les fêtes du printemps]
Le 17 mars, New York fête la Saint-Patrick : sa parade, organisée depuis 1762, est l'une des plus
anciennes du monde. Puis viennent \en{Mother's Day} (deuxième dimanche de mai) et \en{Father's Day}
(troisième dimanche de juin).
\end{culture}
\end{paire}

%% ---------------------------------------------------------------------
\rubrique{pratique}{Exercices}

\exercice[ecrire,niveau=1]{La fête}{Écrivez l'expression anglaise.}
\begin{questions}[2][french]
\item un gâteau → \enb{\trou[34mm]{a cake}}
\item une bougie → \enb{\trou[34mm]{a candle}}
\item faire un vœu → \enb{\trou[34mm]{make a wish}}
\item un cadeau → \enb{\trou[34mm]{a present}}
\item une carte → \enb{\trou[34mm]{a card}}
\item fêter → \enb{\trou[34mm]{celebrate}}
\item Félicitations ! → \enb{\trou[34mm]{Congratulations!}}
\item Santé ! → \enb{\trou[34mm]{Cheers!}}
\end{questions}

\exercice[ecrire,niveau=2]{Le bon temps}{Mettez le verbe au temps qui convient. Aidez-vous des mots
repères.}
\begin{questions}
\item Léa \trou[42mm]{works} at Bright Ideas. \hfill (\en{work})
\item Look! Sam \trou[42mm]{is making} a cake. \hfill (\en{make})
\item Yesterday, Emma \trou[42mm]{bought} a present. \hfill (\en{buy})
\item Next summer, they \trou[42mm]{are going to drive} on Route 66. \hfill (\en{drive})
\item Six months ago, Léa \trou[42mm]{was} very nervous. \hfill (\en{be})
\item \trou[42mm]{Did you see} the candles? \hfill (\en{you / see})
\item Jack \trou[42mm]{doesn't like} coffee. \hfill (\en{not like})
\item Shh! She \trou[42mm]{is coming}! \hfill (\en{come})
\end{questions}

\exercice[ecrire,niveau=2]{Trouvez l'auxiliaire}{Complétez chaque question avec \en{Do}, \en{Does},
\en{Is}, \en{Are}, \en{Can}, \en{Was}, \en{Were} ou \en{Did}.}
\begin{anglais}
\begin{questions}[2]
\item \trou[16mm]{Did} you like the party?
\item \trou[16mm]{Can} you swim?
\item \trou[16mm]{Are} you going to travel?
\item \trou[16mm]{Was} the cake good?
\item \trou[16mm]{Does} Sam work at the café?
\item \trou[16mm]{Is} Léa opening her presents?
\item \trou[16mm]{Were} they happy?
\item \trou[16mm]{Do} you speak French?
\end{questions}
\end{anglais}

\exercice[lire,niveau=2]{Une carte pour Léa}{Mr.\ Carter a glissé une carte dans le cadeau de Léa.
Lisez-la, puis indiquez si chaque phrase est vraie (V) ou fausse (F).}
\begin{anglais}
\otCarteAnniv{Dear Léa,\par\vspace{1mm}
Six months ago, you arrived at Bright Ideas. You didn't know anybody, and your English was a little
shy! Now you are a wonderful designer, and everybody loves your work. Next month, you're going to work
on our biggest project: the new Lumen Books website. Congratulations, and enjoy your day!\par\vspace{1.4mm}
\hfill Ben Carter}
\par\vspace{1mm}\noindent{\fontsize{8.6}{10}\selectfont\color{otMuted}\foreignlanguage{french}{\textit{shy} :
timide · \textit{our biggest project} : notre plus gros projet · \textit{enjoy your day} : profitez bien de
votre journée}\par}
\begin{questions}
\item The card is from Mr.\ Carter. \hfill\trou[9mm]{V}
\item Léa arrived at Bright Ideas a year ago. \hfill\trou[9mm]{F}
\item At first, Léa didn't know anybody. \hfill\trou[9mm]{V}
\item Nobody likes Léa's work. \hfill\trou[9mm]{F}
\item Next month, Léa is going to work on a website. \hfill\trou[9mm]{V}
\item The website is for Lumen Books. \hfill\trou[9mm]{V}
\end{questions}
\end{anglais}

\exercice[ecrire,niveau=2]{Posez la question}{Voici les réponses de Léa. Écrivez la question.}
\begin{questions}
\item I'm twenty-eight. \hfill \enb{\trou[70mm]{How old are you?}}
\item I was born in Lyon. \hfill \enb{\trou[70mm]{Where were you born?}}
\item We went to Boston. \hfill \enb{\trou[70mm]{Where did you go?}}
\item I'm making a cake. \hfill \enb{\trou[70mm]{What are you doing?}}
\item I get up at seven. \hfill \enb{\trou[70mm]{What time do you get up?}}
\item Yes, I can. I can swim very well. \hfill \enb{\trou[70mm]{Can you swim?}}
\end{questions}

\exercice[ecrire,niveau=2]{Corrigez l'erreur}{Chaque phrase contient une erreur. Récrivez-la
correctement.}
\begin{questions}
\item I am work at Bright Ideas. \hfill \enb{\trou[66mm]{I work at Bright Ideas.}}
\item She don't like cake. \hfill \enb{\trou[66mm]{She doesn't like cake.}}
\item Did you went to the party? \hfill \enb{\trou[66mm]{Did you go to the party?}}
\item I'm going travel next year. \hfill \enb{\trou[66mm]{I'm going to travel next year.}}
\item Yesterday I am at home. \hfill \enb{\trou[66mm]{Yesterday I was at home.}}
\end{questions}

\exercice[ecouter,niveau=1,piste=20.10]{Passé, présent ou futur ?}{Écoutez chaque phrase, dite deux
fois. Cochez le temps que vous entendez.}
\begin{questions}[2]
\item \choix{passé}\quad\choix{présent}\quad\choix{futur}
\item \choix{passé}\quad\choix{présent}\quad\choix{futur}
\item \choix{passé}\quad\choix{présent}\quad\choix{futur}
\item \choix{passé}\quad\choix{présent}\quad\choix{futur}
\item \choix{passé}\quad\choix{présent}\quad\choix{futur}
\item \choix{passé}\quad\choix{présent}\quad\choix{futur}
\end{questions}

\exercice[ecouter,ecrire,niveau=2,piste=20.11]{Le discours de Léa}{Léa remercie ses amis. Écoutez et
complétez.}
\begin{anglais}
\begin{questions}
\item Léa arrived in New York six \trou[30mm]{months} ago.
\item At first, she was very \trou[30mm]{nervous}.
\item Now she has wonderful \trou[30mm]{friends}.
\item She loves her \trou[30mm]{job} and she loves New York.
\item Next year, her brother \trou[30mm]{Hugo} is going to visit her.
\end{questions}
\end{anglais}

\exercice[ecouter,ecrire,niveau=2,piste=20.12]{Dictée à trous}{Mrs.\ Lopez parle avec Léa à la fin de
la fête. Écoutez et complétez.}
\begin{anglais}
\noindent\begin{tabular}{@{}l@{\hspace{6mm}}l@{}}
\otNomDit{lopez} & \trou{Did} you have a nice birthday?\\[1.8mm]
\otNomDit{lea}   & Yes! It \trou{was} the best party ever!\\[1.8mm]
\otNomDit{lopez} & What \trou{are} you going to do tomorrow?\\[1.8mm]
\otNomDit{lea}   & I'm \trou{going} to sleep until noon!\\[1.8mm]
\otNomDit{lopez} & Ha ha! \trou{Do} you want some cake for breakfast?\\[1.8mm]
\otNomDit{lea}   & Yes, \trou{please}! Sam \trou{made} a huge cake!\\[1.8mm]
\otNomDit{lopez} & Good night, Léa. And happy \trou{birthday}!\\
\end{tabular}
\end{anglais}

\exercice[ecrire,niveau=2]{Traduisez}{Traduisez en anglais.}
\begin{questions}[1][french]
\item Joyeux anniversaire ! Quel âge as-tu ?\lignes{1}
\item Il y a six mois, je ne parlais pas anglais.\lignes{1}
\item Maintenant, je parle anglais tous les jours.\lignes{1}
\item Regarde ! Ils apportent le gâteau !\lignes{1}
\item L'année prochaine, je vais aller à New York.\lignes{1}
\item Tu as aimé la fête ? — Oui. \textit{(réponse courte)}\lignes{1}
\item Je suis fier de toi !\lignes{1}
\item Merci pour tout !\lignes{1}
\end{questions}

\exercice[ecrire,niveau=3]{Mon bilan}{En six phrases, faites votre propre bilan : votre anglais il y a
quelques mois, ce que vous savez faire maintenant, et ce que vous allez faire ensuite.}
\lignes{6}

\exercice[ecouter,parler,niveau=3,piste=20.13]{À vous de parler !}{Pour terminer ce tome, Emma vous
pose des questions à tous les temps. Répondez \textbf{à voix haute} pendant chaque pause.}
\begin{center}\otTableau
\begin{tabular}{@{}>{\otSemi\fontsize{8.6}{10}\selectfont\color{otNavy}}l@{\hspace{3mm}}>{\enOn}p{70mm}<{\enOff}@{\hspace{6mm}}p{78mm}@{}}
 & \otTH{Emma dit\ldots} & \otTH{Vous\ldots}\\\midrule
1. & What's your name, and where are you from? & présentez-vous.\\
2. & What do you do every day? & décrivez une habitude.\\
3. & What did you do last weekend? & racontez au prétérit.\\
4. & What are you doing right now? & répondez au présent en \en{-ing}.\\
5. & What are you going to do next year? & présentez un projet.\\\bottomrule
\end{tabular}\end{center}

%% ---------------------------------------------------------------------
\rubrique{bilan}{Faites le point}

\begin{essentiel}
\begin{minipage}[t]{0.54\linewidth}
\otEssTitre{Les phrases clés}
\renewcommand{\arraystretch}{1.32}
\begin{tabular}{@{}>{\enOn\otSemi\fontsize{9.4}{11}\selectfont}l<{\enOff}@{\hspace{3mm}}>{\fontsize{9.2}{11}\selectfont\color{white!78!otNavy}}l@{}}
Happy birthday!            & Joyeux anniversaire !\\
Make a wish!               & Fais un vœu !\\
I'm proud of you!          & Je suis fier de toi !\\
Six months ago, I was…     & Il y a six mois, j'étais…\\
Now I speak English.       & Maintenant, je parle anglais.\\
Thank you for everything!  & Merci pour tout !\\
\end{tabular}
\end{minipage}\hfill
\begin{minipage}[t]{0.41\linewidth}
\otEssTitre{Les temps}
{\fontsize{9.2}{11.6}\selectfont \en{I work} · \en{I'm working} · \en{I'm going to work} · \en{I was}
· \en{I worked}, \en{I went}\par}
\otEssTitre{Les questions}
{\fontsize{9.2}{11.6}\selectfont auxiliaire + sujet : \en{Do…? Are…? Can…? Was…? Did…?}\par}
\otEssTitre{Prononciation}
{\fontsize{9.2}{11.6}\selectfont un temps fort par mot important ; les petits mots se compriment\par}
\end{minipage}
\end{essentiel}

\begin{jesais}
\item parler de moi au passé, au présent et au futur ;
\item choisir le bon temps grâce aux mots repères ;
\item poser une question avec le bon auxiliaire et y répondre ;
\item souhaiter un anniversaire, féliciter et remercier ;
\item respecter le rythme de la phrase anglaise ;
\item présenter les grandes fêtes américaines.
\end{jesais}

\clearpage
\begin{testexpress}
\qcm[b]{« J'ai vingt-huit ans » se dit :}{\en{I have twenty-eight years.}}{\en{I'm twenty-eight.}}{\en{I have twenty-eight.}}
\qcm[a]{\en{Look! Sam ……… a cake.}}{\en{is making}}{\en{makes}}{\en{made}}
\qcm[c]{\en{Yesterday, Emma ……… a present.}}{\en{buys}}{\en{is buying}}{\en{bought}}
\qcm[a]{\en{Next summer, I ……… on Route 66.}}{\en{'m going to drive}}{\en{drove}}{\en{drive}}
\qcm[b]{\en{……… you like the party? — Yes, I did.}}{\en{Do}}{\en{Did}}{\en{Were}}
\qcm[c]{\en{……… she work here? — Yes, she does.}}{\en{Is}}{\en{Do}}{\en{Does}}
\qcm[a]{« Fais un vœu ! » se dit :}{\en{Make a wish!}}{\en{Do a wish!}}{\en{Make a vow!}}
\qcm[b]{« Félicitations ! » se dit :}{\en{Cheers!}}{\en{Congratulations!}}{\en{Best regards!}}
\qcm[c]{Dans une phrase anglaise, le rythme est donné par :}{chaque syllabe}{les petits mots}{les syllabes accentuées}
\qcm[a]{Aux États-Unis, on ouvre un cadeau :}{tout de suite, devant la personne}{le lendemain}{seul, en secret}
\end{testexpress}

\noindent{\fontsize{9.6}{12}\selectfont\color{otMuted}Niveau visé (CECRL, niveau A1 atteint !) : « peut
parler de soi au passé, au présent et au futur en phrases simples ». Les corrigés de l'unité sont
page~\pageref{cor-u20-1} ; les textes des pistes d'exercices, page~\pageref{trans-20.10}.\par}
